% Options for packages loaded elsewhere
% Options for packages loaded elsewhere
\PassOptionsToPackage{unicode}{hyperref}
\PassOptionsToPackage{hyphens}{url}
%
\documentclass[
  ignorenonframetext,
]{beamer}
\newif\ifbibliography
\usepackage{pgfpages}
\setbeamertemplate{caption}[numbered]
\setbeamertemplate{caption label separator}{: }
\setbeamercolor{caption name}{fg=normal text.fg}
\beamertemplatenavigationsymbolsempty
% remove section numbering
\setbeamertemplate{part page}{
  \centering
  \begin{beamercolorbox}[sep=16pt,center]{part title}
    \usebeamerfont{part title}\insertpart\par
  \end{beamercolorbox}
}
\setbeamertemplate{section page}{
  \centering
  \begin{beamercolorbox}[sep=12pt,center]{section title}
    \usebeamerfont{section title}\insertsection\par
  \end{beamercolorbox}
}
\setbeamertemplate{subsection page}{
  \centering
  \begin{beamercolorbox}[sep=8pt,center]{subsection title}
    \usebeamerfont{subsection title}\insertsubsection\par
  \end{beamercolorbox}
}
% Prevent slide breaks in the middle of a paragraph
\widowpenalties 1 10000
\raggedbottom
\AtBeginPart{
  \frame{\partpage}
}
\AtBeginSection{
  \ifbibliography
  \else
    \frame{\sectionpage}
  \fi
}
\AtBeginSubsection{
  \frame{\subsectionpage}
}
\usepackage{iftex}
\ifPDFTeX
  \usepackage[T1]{fontenc}
  \usepackage[utf8]{inputenc}
  \usepackage{textcomp} % provide euro and other symbols
\else % if luatex or xetex
  \usepackage{unicode-math} % this also loads fontspec
  \defaultfontfeatures{Scale=MatchLowercase}
  \defaultfontfeatures[\rmfamily]{Ligatures=TeX,Scale=1}
\fi
\usepackage{lmodern}

\usetheme[]{metropolis}
\ifPDFTeX\else
  % xetex/luatex font selection
\fi
% Use upquote if available, for straight quotes in verbatim environments
\IfFileExists{upquote.sty}{\usepackage{upquote}}{}
\IfFileExists{microtype.sty}{% use microtype if available
  \usepackage[]{microtype}
  \UseMicrotypeSet[protrusion]{basicmath} % disable protrusion for tt fonts
}{}
\makeatletter
\@ifundefined{KOMAClassName}{% if non-KOMA class
  \IfFileExists{parskip.sty}{%
    \usepackage{parskip}
  }{% else
    \setlength{\parindent}{0pt}
    \setlength{\parskip}{6pt plus 2pt minus 1pt}}
}{% if KOMA class
  \KOMAoptions{parskip=half}}
\makeatother


\usepackage{longtable,booktabs,array}
\usepackage{calc} % for calculating minipage widths
\usepackage{caption}
% Make caption package work with longtable
\makeatletter
\def\fnum@table{\tablename~\thetable}
\makeatother
\usepackage{graphicx}
\makeatletter
\newsavebox\pandoc@box
\newcommand*\pandocbounded[1]{% scales image to fit in text height/width
  \sbox\pandoc@box{#1}%
  \Gscale@div\@tempa{\textheight}{\dimexpr\ht\pandoc@box+\dp\pandoc@box\relax}%
  \Gscale@div\@tempb{\linewidth}{\wd\pandoc@box}%
  \ifdim\@tempb\p@<\@tempa\p@\let\@tempa\@tempb\fi% select the smaller of both
  \ifdim\@tempa\p@<\p@\scalebox{\@tempa}{\usebox\pandoc@box}%
  \else\usebox{\pandoc@box}%
  \fi%
}
% Set default figure placement to htbp
\def\fps@figure{htbp}
\makeatother





\setlength{\emergencystretch}{3em} % prevent overfull lines

\providecommand{\tightlist}{%
  \setlength{\itemsep}{0pt}\setlength{\parskip}{0pt}}



 


\setbeamerfont{title}{size=\large}
\setbeamertemplate{footline}[frame number]
\makeatletter
\@ifpackageloaded{caption}{}{\usepackage{caption}}
\AtBeginDocument{%
\ifdefined\contentsname
  \renewcommand*\contentsname{Table of contents}
\else
  \newcommand\contentsname{Table of contents}
\fi
\ifdefined\listfigurename
  \renewcommand*\listfigurename{List of Figures}
\else
  \newcommand\listfigurename{List of Figures}
\fi
\ifdefined\listtablename
  \renewcommand*\listtablename{List of Tables}
\else
  \newcommand\listtablename{List of Tables}
\fi
\ifdefined\figurename
  \renewcommand*\figurename{Figure}
\else
  \newcommand\figurename{Figure}
\fi
\ifdefined\tablename
  \renewcommand*\tablename{Table}
\else
  \newcommand\tablename{Table}
\fi
}
\@ifpackageloaded{float}{}{\usepackage{float}}
\floatstyle{ruled}
\@ifundefined{c@chapter}{\newfloat{codelisting}{h}{lop}}{\newfloat{codelisting}{h}{lop}[chapter]}
\floatname{codelisting}{Listing}
\newcommand*\listoflistings{\listof{codelisting}{List of Listings}}
\makeatother
\makeatletter
\makeatother
\makeatletter
\@ifpackageloaded{caption}{}{\usepackage{caption}}
\@ifpackageloaded{subcaption}{}{\usepackage{subcaption}}
\makeatother

\usepackage{bookmark}
\IfFileExists{xurl.sty}{\usepackage{xurl}}{} % add URL line breaks if available
\urlstyle{same}
\hypersetup{
  hidelinks,
  pdfcreator={LaTeX via pandoc}}


\author{}
\date{}

\begin{document}


\begin{frame}
% TITLE SLIDES
\title{\Large La mia esperienza con i dati FAIR}
\author{Giorgio Arcara}

\titlegraphic{
\vspace*{7cm}
\includegraphics[scale=0.15]{Figures/LogoSanCamilloIRCCS_Unipd_alpha.png}
\begin{center}
\vspace{-2.2em}
\includegraphics[scale=0.15]{Figures/CC_license_3_0.png}
\end{center}
}
\date{21/05/2026}
\maketitle
\end{frame}

\begin{frame}{Due esperienze}
\phantomsection\label{due-esperienze}
\begin{itemize}
\tightlist
\item
  Ricercatore
\item
  Vice Direttore Scientifico
\end{itemize}
\end{frame}

\section{Ricercatore}\label{ricercatore}

\begin{frame}{Utilizzatore di dati FAIR}
\phantomsection\label{utilizzatore-di-dati-fair}
\pause

Serve veramente raccogliere dei nuovi dati?
\end{frame}

\begin{frame}{Un primo esempio}
\phantomsection\label{un-primo-esempio}
\textbf{Aperiodic component} of EEG spectrum in \textbf{Aging} \pause
\pandocbounded{\includegraphics[keepaspectratio]{Figures/montemurro_aperiodic.png}}
\centering
\href{https://www.nature.com/articles/s41598-024-66049-2}{\underline{Montemurro et al., 2024}}
\pause

\href{https://www.nature.com/articles/sdata2018308}{\underline{\textbf{Lemon dataset}}}
\end{frame}

\begin{frame}{Un primo esempio - risultati}
\phantomsection\label{un-primo-esempio---risultati}
\begin{center}
\includegraphics[width=0.8\linewidth,height=\textheight,keepaspectratio]{Figures/aperiodic_montemurro.png}
\end{center}
\end{frame}

\begin{frame}{Alcuni esempi}
\phantomsection\label{alcuni-esempi}
\pandocbounded{\includegraphics[keepaspectratio]{Figures/Lago_aperiodic.png}}
\centering
\href{https://journals.plos.org/plosone/article?id=10.1371/journal.pone.0328318}{\underline{Lago et al., 2026}}
\pause

\href{https://github.com/euroladbrainlat/Brain-health-in-diverse-setting}{\underline{\textbf{Brain health in diverse setting - dataset}}}
\end{frame}

\begin{frame}{Un altro esempio}
\phantomsection\label{un-altro-esempio}
\pandocbounded{\includegraphics[keepaspectratio]{Figures/Schelfhout_fingerprinting.png}}
\centering
\href{https://www.biorxiv.org/c/10.1101/2024.08.05.606621v1}{\underline{Schelfhout et al., 2026}}

\pause

\href{https://cam-can.mrc-cbu.cam.ac.uk/dataset/}{\underline{\textbf{Cam-Can data repository}}}
\href{https://www.mcgill.ca/bic/neuroinformatics/omega}{\underline{\textbf{Omega Dataset}}}
\end{frame}

\begin{frame}{E un altro ancora}
\phantomsection\label{e-un-altro-ancora}
A PhD on DBS data
\href{https://figshare.com/articles/dataset/Sumultaneous_subthalamic_local_field_potentials_and_MEG-derived_cortical_activity_in_Parkinson_s_patients_ON_and_OFF_dopaminergic_medication/31254136}{\underline{\textbf{DBS - ON/OFF Dataset}}}
\end{frame}

\begin{frame}{A good habit}
\phantomsection\label{a-good-habit}
L'atto della Fairification non è \textbf{solo} importante per l' Open
Science e per aspetti ideologici.\pause

\pause

E' un ``good habit'' che fa \emph{risparmiare tempo}. \pause

\begin{itemize}
\tightlist
\item
  Non vengono raccolti dati non necessari
\item
  porta ad un migliore (e più sicuro) immagazzinamento dei dati, anche
  per riutilizzo interno.
\end{itemize}
\end{frame}

\section{Vice direttore scientifico}\label{vice-direttore-scientifico}

\begin{frame}{Perchè FAIR data?}
\phantomsection\label{perchuxe8-fair-data}
Quando si vuole parlare di \emph{impatto} sempre più sono necessari
grandi moli di dati. I dati fair sono un pre-requisito \pause

Avere dati FAIR significa anche essere più pronti a collaborazioni.
\end{frame}

\begin{frame}{FAIR data, Data Management Plan}
\phantomsection\label{fair-data-data-management-plan}
Numerosi finanziamenti, richiedono un DMP e dati FAIR.

Conoscere pratiche FAIR significa stare al passo con i tempi
\end{frame}

\begin{frame}{Grazie!}
\phantomsection\label{grazie}
\centering
\Large

giorgio.arcara@unipd.it

\url{giorgioarcara.github.io}

\vfill
\begin{columns}
\column{0.5\textwidth}
\begin{center}
\includegraphics[width=0.5\textwidth]{Figures/ONDA_logo_v12a.png}

\href{https://sancamilloirccs-ondalab.github.io/}{\underline{\textbf{ONDA Lab}}}
\end{center}
\column{0.5\textwidth}
\begin{center}
\includegraphics[width=0.3\textwidth]{Figures/QP_logo.png}

\href{https://sancamilloirccs-ondalab.github.io/}{\underline{\textbf{QP Lab}}}
\end{center}
\end{columns}
\end{frame}




\end{document}
